\section{Cấu Trúc Báo Cáo}

Báo cáo được tổ chức thành các chương với nội dung và vai trò rõ ràng nhằm trình bày tuần tự từ bối cảnh bài toán đến thiết kế, triển khai và đánh giá hệ thống:
\begin{itemize}
    \item \textbf{Chương 1 – Giới thiệu}: Trình bày bối cảnh bài toán, động lực thực hiện, mục tiêu của đồ án, phạm vi và cấu trúc tổng thể của báo cáo.
    \item \textbf{Chương 2 – Các công trình liên quan}: Tổng quan các hướng tiếp cận và hệ thống tiêu biểu trong lĩnh vực xử lý hội thoại nhiều người, nhận dạng tiếng nói và truy xuất thông tin, làm cơ sở định vị hướng tiếp cận của đồ án.
    \item \textbf{Chương 3 – Cơ sở lý thuyết}: Trình bày các khái niệm và nền tảng kỹ thuật cần thiết như phát hiện hoạt động giọng nói, biểu diễn embedding, phân cụm, nhận dạng tiếng nói và truy xuất văn bản, phục vụ cho việc hiểu các chương sau.
    \item \textbf{Chương 4 – Thiết kế và hiện thực hệ thống}: Mô tả kiến trúc tổng thể của hệ thống, luồng dữ liệu và thiết kế chi tiết các mô-đun chính, bao gồm mô-đun xử lý tiếng nói và mô-đun truy xuất và tổng hợp thông tin.
    \item \textbf{Chương 5 – Đánh giá và thảo luận}: Trình bày các kịch bản đánh giá, tiêu chí đo lường và kết quả thực nghiệm nhằm so sánh các cấu hình hệ thống và phân tích các yếu tố ảnh hưởng đến hiệu năng.
    \item \textbf{Chương 6 – Kết luận}: Tổng kết các kết quả đạt được, nêu hạn chế của hệ thống và định hướng phát triển trong tương lai.
\end{itemize}
